\thispagestyle{empty}
\section*{Kurzdarstellung}
\label{sec:kurzdarstellung}

Dieser Arbeit entwickelte ein statistisches Modell, um die geometrischen Eigenschaften von Cookies zu prognostizieren.
Aufgrund der Limitierungen der Arbeit wird der Cookie-Durchmesser als Zielgröße anhand der Faktoren Zucker, Butter und Vollkornmehlanteil analysiert.
Der Versuchsplan ist ein 2³-Vollfaktorplan mit Zentralpunkt und drei Blöcken für die Replikationen.
Ein Backvorgang enthält die neun FSK des Versuchsplans und stellt eine Replikation dar.
Die Positionen der FSK auf dem Backblech werden erst randomisiert und dann rotiert um die ungleiche Wärmeentwicklung im Backofen ausgleichen.
Allerdings spiegeln die Residuen des Modells diese Ungleichheit durch Ausreißer wieder.
Über die Vereinfachung wird das Modell auf die drei Haupteffekten und Blöcke reduziert.
Das Ergebnis ist ein signifikantes Modell für das Signifikanzniveau von 1\,\%. 
Das Modell erzielt einen R²-Wert von 96,55\,\% und einen korrigierten R²-Wert von 95,83\,\%.
%Einen Hinweis auf Krümmung gab es beim Niveau von 1\,\% nicht.